\section{Open problems}\label{sec:open-problems}
The proof uses the coefficient distribution only through (H1)--(H5),
bounds its errors by low moments, and sees the zeros only at the scale
$1/n$. This suggests three problems: the weakest moment condition, the
size of the fluctuations, and the local statistics of the zeros.

\begin{problem}[The moment threshold]\label{prob:threshold}
Let $(\xi_k)_{k\ge0}$ be independent and identically distributed complex
random variables, not almost surely constant, and let $f_n$ and $\nu_n$ be
as in Theorem~\ref{thm:radial-profile}. Does the conclusion of
Theorem~\ref{thm:radial-profile} hold whenever $\mathbb E|\xi_0|^2<\infty$?
\end{problem}

Theorem~\ref{thm:radial-profile} establishes the law for five families of
coefficient distributions, and our proof uses the distribution only through
the hypotheses (H1)--(H5) of Section~\ref{sec:layer2-models}, which hold
for signs by Sections~\ref{sec:concentration} and~\ref{sec:local-roots}
and which Proposition~\ref{prop:coefficient-transfer} verifies for the
other four families.
The second moment cannot be weakened.
If $\xi_n\ne0$, Jensen's formula for $w^nf_n(1/w)$ on the circle
$|w|=e^{1/n}$ and the concavity of the logarithm give
\[
 n-\nu_n(1)\le\frac n2\log\Bigl(1+e^2\sum_{k<n}\frac{|\xi_k|^2}{|\xi_n|^2}\Bigr).
\]
When $\mathbb E|\xi_0|^2=\infty$, the ratio $|\xi_n|^2/\sum_{k<n}|\xi_k|^2$
is unbounded almost surely, by a theorem of Kesten~\cite{Kesten1970}, and
hence $\limsup_n\nu_n(1)/n=1$.
At a fixed degree, with mean zero and a finite second moment, the zeros
near each point of the unit circle converge in distribution at the scale
$1/n$~\cite{KabluchkoKhoruzhenkoMarynych2025}, and equidistribution of the
zeros on the unit circle needs only a logarithmic
moment~\cite{IbragimovZaporozhets2013}.

\begin{problem}[Fluctuations]\label{prob:fluctuations}
For Rademacher coefficients, is $(\nu_n(1)-n/2)/\sqrt n$ asymptotically
normal with a positive variance? Is
\[
 \limsup_{n\to\infty}\frac{\nu_n(1)-n/2}{\sqrt{n\log\log n}}
\]
a positive and finite constant almost surely?
\end{problem}

Theorem~\ref{thm:log-rate} gives $\nu_n(1)-n/2=O(n/\log n)$ almost surely.
Positive answers would identify the exact almost-sure order
$\sqrt{n\log\log n}$, as in the law of the iterated logarithm for sums of
independent random variables.
Along the nested sequence, the increments $\nu_{n+m}(1)-\nu_n(1)$ measure
the flow of zeros through the unit circle as coefficients are appended,
and their size is part of the question.

\begin{problem}[Almost-sure local universality]\label{prob:local}
Let the coefficients be as in Theorem~\ref{thm:radial-profile}, with
$\mathbb E|\xi_0|^2=1$, and let $G(\zeta)=\int_0^1e^{t\zeta}\dd W(t)$,
where $W$ is a complex Brownian motion with $\mathbb E|W(t)|^2=t$.
Almost surely, do the correlation functions of the zeros of
$\zeta\mapsto f_n(e^{i\theta}(1+\zeta/n))$, averaged over $\theta$,
converge to those of the zeros of $G$?
\end{problem}

The zeros of $G$ have intensity $\Phi'(\operatorname{Re}\zeta)/(2\pi)$, and
Theorem~\ref{thm:radial-profile} is the one-point case of the problem.
At a fixed degree, local universality of the zeros near the unit circle
was proved by Tao and Vu~\cite{TaoVu2015}, and in distribution under a
finite second moment by Kabluchko, Khoruzhenko and
Marynych~\cite{KabluchkoKhoruzhenkoMarynych2025}. The problem asks for it
almost surely, at every degree of one nested sequence.
In the proof of Theorem~\ref{thm:radial-profile}, root matching
(Theorem~\ref{thm:T4}) keeps each selected zero of $f_N$ near the unit
circle within $N^{-65/64}=o(1/N)$ of a zero of $f_n$, at every degree $n$
of its block, so that local configurations at the sparse degrees persist
at every degree.
